\section{Open Problems}
\label{sec:open_problems}

This section catalogues the principal open problems identified in the review,
organised by theme.
For each problem we note the strongest evidence motivating it and the specific
gap in current knowledge.

\paragraph{Priority ordering}
The nine problems are not of equal urgency, and urgency must be distinguished from blocking: these are different properties operating on different timescales.
RP1 (standardised evaluation) is the most broadly \emph{blocking}: without a shared protocol, no other open problem can be resolved in a reproducible way, because progress claimed under one evaluation setup cannot be compared to progress under another.
MM3 (audio and video safety) is the most \emph{urgent}: capable multimodal models with audio and video capabilities are already in production deployment --- GPT-4o audio mode, Gemini Live --- with no academic safety characterisation at all.
The tension between these two rankings is only apparent.
RP1 blocking and MM3 urgency operate on different axes: RP1 is a methodological prerequisite for \emph{measuring and comparing} progress across future work, while MM3 describes the gap between deployed-system risk and current academic coverage.
MM3 is urgent not because the problem is well-posed and tractable (it is neither), but because, unlike streaming text safety where the problem is at least defined and the attack surface understood, audio and video safety lacks even a foundational attack-surface taxonomy; the distance between deployment reality and academic treatment is widest precisely here.
MM3 urgency therefore calls for immediate investment in problem definition --- an attack-surface taxonomy for audio-language and video-language models --- even before a reproducible evaluation protocol for that domain exists.
RP1 and MM3 are thus complementary rather than competing priorities: MM3 motivates \emph{why} we need the field to move fast, and RP1 explains \emph{how} that movement must be structured to produce reproducible results.
MM1--MM2 (cross-modal alignment and multi-turn VLM attacks) require more foundational work on VLM architecture before targeted guardrail solutions are feasible.
FT1 (fine-tuning robustness) and FG1 (formal guarantees) are the most foundationally difficult: FT1 requires solving an optimisation-theoretic problem about weight-space invariants, and FG1 faces the non-stationary adversary obstacle described below.
RT1--RT3 (streaming safety) are the most tractable in the near term: they are empirical questions that could be addressed by a single rigorous study, and the adjacent constrained-generation literature provides structural inspiration.
They are also \emph{sequentially dependent}: RT1 (minimum viable prefix length) is a prerequisite for RT2 and RT3 --- without knowing the minimum prefix at which a classifier can operate reliably, the parallelism architecture (RT2) has no design target, and the early-stopping cost analysis (RT3) has no threshold to optimise against.
Future work should therefore address RT1 before RT2 or RT3.

\subsection{The Reproducibility Crisis}
\label{subsec:reproducibility}

The most pervasive finding of this review is methodological.
Across three temporally separated author-verification passes spanning 75 primary
sources and 237 extracted claims, exactly five claims survived structured self-review
(unanimous agreement across all three passes required).

Table~\ref{tab:claim_counts} breaks down the 237 extracted claims by section to
allow verification that the section-level totals are consistent with the headline
count.

\begin{table}[h]
  \caption{Section-level breakdown of the 237 extracted claims.
           Claims are quantitative or qualitative architectural assertions only;
           purely descriptive or background statements were not coded as discrete claims.
           The fifth surviving claim (evaluation-protocol fragility, SALAD-Bench
           meta-finding) is a cross-cutting meta-finding drawn from the 72 remaining
           papers; it is enumerated in the text below and in Section~\ref{sec:conclusion}
           but is not attributed to a single row in this table.}
  \label{tab:claim_counts}
  \centering
  \begin{tabular}{lrr}
    \toprule
    \textbf{Section} & \textbf{Claims extracted} & \textbf{Claims surviving} \\
    \midrule
    \S3 Background and Taxonomy         & 18  & 2 \\
    \S4 Real-Time Guardrails            & 25  & 1 \\
    \S5 Multimodal Guardrails           & 26  & 2 \\
    \S6 Reproducibility (meta-finding)  &  0  & 0 \\
    Remaining 72 corpus papers          & 168 & 0 \\
    \midrule
    \textbf{Total}                      & \textbf{237} & \textbf{5} \\
    \bottomrule
  \end{tabular}
  \par\smallskip
  \noindent\small The 168 remaining claims are distributed across the 72 corpus
  papers not individually highlighted in Sections~\ref{sec:background}--\ref{sec:multimodal}.
  The SALAD-Bench meta-finding (fifth surviving claim) is the only surviving claim
  from this group; it is reported separately because it is a cross-paper methodological
  observation, not a result of any single paper.
\end{table}

The five surviving claims are:

\begin{enumerate}
  \item The dual reward model architecture and threshold arbitration in Llama 2
        ($R_s < 0.15$)~\cite{touvron2023llama2}, verified in
        Section~\ref{sec:realtime} (Dual Reward Model Architecture subsection).
  \item Safety Context Distillation as a weight-embedding technique~\cite{touvron2023llama2},
        verified in Section~\ref{sec:background} (Foundational Alignment Techniques
        subsection).
  \item The Llama Guard instruction-following classification
        architecture~\cite{inan2023llama}, verified in Section~\ref{sec:realtime}
        (Llama Guard subsection).
  \item Multi-turn conversational priming as a jailbreak amplifier (+29--71\% over
        single-turn attacks)~\cite{russinovich2024greatescalation,li2024llm},
        verified in Section~\ref{sec:multimodal} (Multi-Turn Conversational Attack
        subsection).
  \item Evaluation methodology choices shifting reported attack-success rates by up
        to 30 percentage points~\cite{li2024salad}, verified as a meta-finding
        across Sections~\ref{sec:realtime} and~\ref{sec:multimodal}.
\end{enumerate}

The pattern of refutations is instructive.
Refuted claims are not uniformly fabricated; they reflect genuine experimental
observations that collapse under three failure modes:
(i) \emph{evaluation-protocol fragility}---at least 9 distinct evaluation setups
appear across papers, and token-output budget alone accounts for $\pm$30\% variation
in ASR~\cite{li2024salad};
(ii) \emph{model-version sensitivity}---results reported for GPT-4-0314 may not hold
for GPT-4-0613 or subsequent updates;
(iii) \emph{overstated precision}---specific numerical claims (e.g.\ exact ASR
percentages) cannot be reproduced from available source material.

\paragraph{Open problem (RP1)}
\textit{The field lacks a standardised, versioned evaluation protocol for runtime
guardrail effectiveness.}
Current practice uses incompatible combinations of attack sets, judge models,
token budgets, and harm taxonomies, making cross-paper comparisons unreliable.
Several recent frameworks attempt to address this: HarmBench~\cite{mazeika2024harmbench}
proposes a standardised red-teaming evaluation with fixed attack sets and a common
judge, finding that attack success rates vary dramatically across evaluation setups;
StrongREJECT~\cite{souly2024strongreject} identifies a specific failure mode of
common jailbreak evaluations---empty or low-quality jailbreaks scoring as successful
when judge models apply lenient criteria; GRAM~\cite{zhu2024gram} analyses
generalisation failures in safety evaluation by showing that guardrails overfit to
training attack distributions; and EasyJailbreak~\cite{zhou2024easyjailbreak} provides
a unified framework exposing how the choice of jailbreak methodology alone accounts
for substantial variance in reported attack success rates.
Collectively these works confirm that the evaluation fragility is not incidental but
structural: progress claimed by any paper that does not fix all five protocol
dimensions (attack set version, judge model version, token budget, harm taxonomy,
and ambiguous-output handling) is not reproducible.
A standardised protocol must specify all five dimensions before any guardrail results
can be compared across papers.

\subsection{The Streaming Safety Gap}

As established in Section~\ref{subsec:streaming_gap}, no verified mechanism exists
for enforcing safety constraints on streamed token sequences in real time.
(This null finding was confirmed across three independent search passes using streaming-specific query terms including ``token-by-token safety,'' ``prefix classification,'' ``streaming guardrail,'' and ``in-generation filtering''; the absence reflects the state of the literature, not a gap in the search strategy.)

\paragraph{Open problem (RT1) \textnormal{[prerequisite for RT2 and RT3]}}
\textit{Minimum viable prefix length for streaming safety detection.}
For each harm category (explicit content, self-harm facilitation, cyberweapon
synthesis, etc.), what is the minimum prefix length at which a classifier can
achieve acceptable precision and recall?
This determines whether in-generation safety checking is architecturally feasible
without unacceptable latency or over-refusal rates.
RT1 is the foundational empirical question: the answer sets the design target for
the parallelism architecture (RT2) and the minimum window for the stopping policy
analysis (RT3).
Progress on RT2 or RT3 without first answering RT1 risks optimising an architecture
with no empirically grounded operating point.

A deployment caveat applies to any empirical RT1 finding: the minimum viable prefix length, once published, becomes an adversarial target.
An attacker who knows the threshold $L^*$ can front-load the first $L^* - 1$ tokens with benign content and inject policy-violating material thereafter.
RT1 answers should therefore be treated as private security parameters rather than published design constants --- analogous to the non-stationarity critique that applies to Aegis's no-regret bound (Section~\ref{sec:open_problems}).
Future work on RT1 should include an adversarial robustness analysis that evaluates classifier performance when the threshold is known to the attacker, not only under the standard oblivious-attacker assumption.

\paragraph{Open problem (RT2) \textnormal{[requires RT1 answer]}}
\textit{Parallelism architecture for in-generation safety checking.}
A streaming-compatible safety mechanism must run concurrently with the generation
process.
Whether this requires a separate inference process, speculative decoding tricks,
or a fundamentally different model architecture is an open design question whose
feasibility depends on the minimum prefix length established by RT1.

\paragraph{Open problem (RT3) \textnormal{[requires RT1 answer]}}
\textit{Early-stopping cost under uncertainty.}
Conservative in-generation classifiers may stop generation on ambiguous prefixes,
trading safety for utility.
The optimal stopping policy under uncertainty---and its sensitivity to harm category,
user population, and deployment context---has not been studied.
The RT1 answer determines the earliest prefix at which an early-stopping decision
is meaningful; without it, the policy optimisation in RT3 has no grounded lower bound.

\subsection{The Cross-Modal Alignment Deficit}

\paragraph{Open problem (MM1)}
\textit{Does safety alignment in the text backbone transfer to visual and audio
modalities?}
The qualitative concern that multimodal fine-tuning introduces new vulnerability
pathways has not been rigorously characterised.
A reproducible measurement protocol---one that isolates the contribution of
multimodal fine-tuning from changes in the training data distribution---is needed.

\paragraph{Open problem (MM2)}
\textit{Multi-turn conversational attacks on VLMs.}
The confirmed finding that multi-turn priming amplifies jailbreaks in text-only
models~\cite{russinovich2024greatescalation} has not been evaluated in VLM settings.
Does image context accelerate or retard multi-turn escalation?
Does visual priming (a sequence of benign images) create the same vulnerability as
text priming?

\paragraph{Open problem (MM3)}
\textit{Audio and video safety mechanisms.}
No verified safety mechanism for audio-language or video-language models exists in
the literature.
Given active deployment of audio-capable models, this represents the most urgent
gap in the multimodal safety landscape.

\subsection{Post-Deployment Fine-Tuning Threats}

\paragraph{Open problem (FT1)}
\textit{Safety alignment robustness under post-deployment customisation.}
Qi et al.~\cite{qi2023finetuning} raise the concern that fine-tuning on entirely
benign data may degrade safety alignment.
The specific numerical claims in that paper did not survive adversarial verification,
but the threat model is sound: any deployment context that allows users to fine-tune
the model (e.g.\ OpenAI fine-tuning API, open-source model adaptation) creates a
potential mechanism for safety degradation.
What architectural invariants would make safety behaviour robust to post-deployment
fine-tuning?

\subsection{Formal Guarantees}

\paragraph{Open problem (FG1)}
\textit{Completeness bounds for runtime guardrails.}
No verified formal result provides bounds on the false-negative rate of any runtime
guardrail against novel jailbreaks outside the training distribution.
The Aegis no-regret bound~\cite{su2024aegis} is the closest attempt but did not
survive adversarial verification.
Establishing even weak distribution-dependent bounds---analogous to PAC-learning
guarantees---would represent a significant theoretical advance.

The most instructive case study is erase-and-check~\cite{kumar2023certifying}, which does
provide formal guarantees --- but under a specific threat model that does not match
operational deployment.
Erase-and-check provides certified guarantees under a bounded random-erasure threat
model: it can certify that no harmful content survives a specified erasure budget.
These guarantees do not, however, extend to adaptive adversaries who craft inputs
specifically to survive erasure with high probability --- through paraphrasing,
insertion, semantic-preserving transformations, or embedding-space perturbations.
This is not a limitation of the technique within its threat model; it is a
\emph{threat-model mismatch} between what erase-and-check certifies (robustness to
bounded random erasure) and what deployed guardrails must withstand (adaptive
adversarial attack).
Practitioners should therefore treat erase-and-check certificates as providing strong
theoretical assurance under restricted assumptions, not as practical robustness
guarantees against motivated adversaries.

Two structural obstacles explain why this threat-model mismatch is not easily
closed by extending the erase-and-check framework: (i) the adversary is
\emph{non-stationary} --- unlike the fixed but unknown distribution assumed in
classical PAC learning, a jailbreak adversary adapts to the deployed guardrail,
so the adaptive attacker who crafts inputs to survive erasure is precisely the
adversary that falls outside the bounded random-erasure model, rendering any
bound derived on a fixed attack distribution invalid after the first adaptive
attack iteration; and (ii) the input distribution over jailbreak prompts is
\emph{undefined by default} --- there is no agreed natural distribution over
harmful queries, making it unclear what an $\epsilon$-$\delta$ bound would be
taken over even in the non-adaptive case.
Future work on FG1 should therefore frame completeness bounds explicitly in
adversarial or minimax terms --- worst-case regret against a computationally
bounded adaptive adversary --- rather than in the classical PAC formulation,
which assumes a fixed and non-adaptive attack distribution.
